% GENERATED FILE. Edit canonical lessons, then run npm run generate:books.
% canonical-source-hash: fnv1a64:eefe88b6915d05ba
% canonical-lessons: GE-C226-aussprache, GE-C226-bedeutung, GE-C226-begruessung, GE-C226-beschreibung, GE-C226-besuch

\chapter{Pronunciation, A meaning, A greeting, A description, A visit}
\label{ch:ge-pronunciation-a-meaning-a-greeting-a-description-a-visit}
\hlchaptermodality{226}

\begin{chapteropening}
\textbf{By the end of this chapter:} \emph{I can say pronunciation, a meaning, a greeting, a description and a visit.}

Complete the last lesson of chapter 226: i can say pronunciation, a meaning, a greeting, a description and a visit.
\end{chapteropening}

\section[die Aussprache]{die Aussprache — pronunciation}

\label{lesson:GE-C226-aussprache}



\begin{quote}
Before the new one: say the German for training, then the German for an exhibition.
\end{quote}

\subsection*{You'll want to know: die Aussprache}
\textbf{die Aussprache} — \textquotedblleft{}pronunciation\textquotedblright{}.

\begin{grammarlens}[title={What you've built}]
One more word for this chapter.
\end{grammarlens}

\subsection*{Your turn}
\begin{itemize}
\raggedright
  \item \emph{Say it:} \emph{die Aussprache}
  \item \emph{Say it:} \emph{die Aussprache}, once more
  \item \emph{Say it:} say \emph{die Aussprache}
  \item \emph{Recall:} say the German for training, then the German for an exhibition, then say \emph{die Aussprache} again
\end{itemize}

\begin{tcolorbox}[breakable,colback=teal!4,colframe=teal!35!black,title={Before you move on}]
What does die Aussprache mean? (\textquotedblleft{}Pronunciation\textquotedblright{}.) Say it once more.
\end{tcolorbox}
\section[die Bedeutung]{die Bedeutung — a meaning}

\label{lesson:GE-C226-bedeutung}



\begin{quote}
Before the new one: say the German for an exhibition, then the German for pronunciation.
\end{quote}

\subsection*{You'll want to know: die Bedeutung}
\textbf{die Bedeutung} — \textquotedblleft{}a meaning\textquotedblright{}.

\begin{grammarlens}[title={What you've built}]
One more word for this chapter.
\end{grammarlens}

\subsection*{Your turn}
\begin{itemize}
\raggedright
  \item \emph{Say it:} \emph{die Bedeutung}
  \item \emph{Say it:} \emph{die Bedeutung}, once more
  \item \emph{Say it:} say \emph{die Bedeutung}
  \item \emph{Recall:} say the German for an exhibition, then the German for pronunciation, then say \emph{die Bedeutung} again
\end{itemize}

\begin{tcolorbox}[breakable,colback=teal!4,colframe=teal!35!black,title={Before you move on}]
What does die Bedeutung mean? (\textquotedblleft{}A meaning\textquotedblright{}.) Say it once more.
\end{tcolorbox}
\section[die Begrüßung]{die Begrüßung — a greeting, a welcome}

\label{lesson:GE-C226-begruessung}



\begin{quote}
Before the new one: say the German for pronunciation, then the German for a meaning.
\end{quote}

\subsection*{You'll want to know: die Begrüßung}
\textbf{die Begrüßung} — \textquotedblleft{}a greeting, a welcome\textquotedblright{}.

\begin{grammarlens}[title={What you've built}]
One more word for this chapter.
\end{grammarlens}

\subsection*{Your turn}
\begin{itemize}
\raggedright
  \item \emph{Say it:} \emph{die Begrüßung}
  \item \emph{Say it:} \emph{die Begrüßung}, once more
  \item \emph{Say it:} say \emph{die Begrüßung}
  \item \emph{Recall:} say the German for pronunciation, then the German for a meaning, then say \emph{die Begrüßung} again
\end{itemize}

\begin{tcolorbox}[breakable,colback=teal!4,colframe=teal!35!black,title={Before you move on}]
What does die Begrüßung mean? (\textquotedblleft{}A greeting, a welcome\textquotedblright{}.) Say it once more.
\end{tcolorbox}
\section[die Beschreibung]{die Beschreibung — a description}

\label{lesson:GE-C226-beschreibung}



\begin{quote}
Before the new one: say the German for a meaning, then the German for a greeting.
\end{quote}

\subsection*{You'll want to know: die Beschreibung}
\textbf{die Beschreibung} — \textquotedblleft{}a description\textquotedblright{}.

\begin{grammarlens}[title={What you've built}]
One more word for this chapter.
\end{grammarlens}

\subsection*{Your turn}
\begin{itemize}
\raggedright
  \item \emph{Say it:} \emph{die Beschreibung}
  \item \emph{Say it:} \emph{die Beschreibung}, once more
  \item \emph{Say it:} say \emph{die Beschreibung}
  \item \emph{Recall:} say the German for a meaning, then the German for a greeting, then say \emph{die Beschreibung} again
\end{itemize}

\begin{tcolorbox}[breakable,colback=teal!4,colframe=teal!35!black,title={Before you move on}]
What does die Beschreibung mean? (\textquotedblleft{}A description\textquotedblright{}.) Say it once more.
\end{tcolorbox}
\section[der Besuch]{der Besuch — a visit}

\label{lesson:GE-C226-besuch}



\begin{quote}
Before the new one: say the German for a greeting, then the German for a description.
\end{quote}

\subsection*{You'll want to know: der Besuch}
\textbf{der Besuch} — \textquotedblleft{}a visit\textquotedblright{}.

\begin{grammarlens}[title={What you've built}]
One more word for this chapter.
\end{grammarlens}

\subsection*{Your turn}
\begin{itemize}
\raggedright
  \item \emph{Say it:} \emph{der Besuch}
  \item \emph{Say it:} \emph{der Besuch}, once more
  \item \emph{Say it:} say \emph{der Besuch}
  \item \emph{Recall:} say the German for a greeting, then the German for a description, then say \emph{der Besuch} again
\end{itemize}

\begin{tcolorbox}[breakable,colback=teal!4,colframe=teal!35!black,title={Before you move on}]
What does der Besuch mean? (\textquotedblleft{}A visit\textquotedblright{}.) Say it once more.
\end{tcolorbox}
